%=====================================================================
\section{Simulation: event processes and scenarios}\label{sec:v71sim}
% numbered from 1820: the neural generators use 1800-1805
\newcounter{savedpropsim}\setcounter{savedpropsim}{\value{proposition}}
\setcounter{proposition}{1819}
This section adds two kinds of simulation to the engine. \emph{Event processes} extend the sequence model of
Section~\ref{sec:v32ts}: a Bayesian Hawkes process for self-exciting histories (\texttt{sequence.process: hawkes},
\texttt{relsynth/timeseries/hawkes.py}) and a multi-state semi-Markov chain for a status column of the events
(\texttt{sequence.states}, \texttt{relsynth/timeseries/states.py}). \emph{Scenarios} (\texttt{synthesis.scenarios},
\texttt{relsynth/scenarios.py}) are spec-driven interventions applied to a posterior draw while a dataset is sampled.
Both are fitted from data or act on the fitted model; nothing is tuned per dataset, and every choice is a spec key.
Guides: \texttt{docs/timeseries.md}, \texttt{docs/scenarios.md}. Tests: \texttt{tests/unit/test\_event\_processes.py},
\texttt{tests/unit/test\_scenarios.py}.

Notation follows Section~\ref{sec:v32ts}: entity $i$ of segment $k$ has $n_i$ events at operational times
$u_{i1}\le\dots\le u_{in_i}$; its first and last events come from the life-time law, $S_i=u_{in_i}-u_{i1}$.

\begin{proposition}[Hawkes likelihood and profile baseline]\label{prop:simhawkesfit}
Let the interior events of entity $i$ on $(u_{i1},u_{in_i})$, with $u_{i1}$ in the history, follow the intensity
$\lambda_i(u)=\mu_i+\eta\beta\sum_{u_{ik}<u}e^{-\beta(u-u_{ik})}$ with $0\le\eta<1$, $\beta>0$. Then
\[
\ell_i(\mu_i,\eta,\beta)=\sum_{1<j<n_i}\log(\mu_i+\eta\beta A_{ij})-\mu_iS_i
-\eta\sum_{j<n_i}\bigl(1-e^{-\beta(u_{in_i}-u_{ij})}\bigr),\qquad
A_{ij}=e^{-\beta(u_{ij}-u_{i,j-1})}(1+A_{i,j-1}),\ A_{i1}=0,
\]
is the log likelihood of the interior events, $\ell_i$ is strictly concave in $\mu_i\ge0$, and its maximiser is the
unique root of $\sum_j(\mu+\eta\beta A_{ij})^{-1}=S_i$ when the left side at $\mu=0$ exceeds $S_i$, and $0$
otherwise. Newton steps in $x=\log\mu$ started at $\max(\bar n/S_i-\bar e,\epsilon)$ (with $\bar n$ the number of
terms and $\bar e$ their mean excitation) converge to it.
\end{proposition}
\begin{proof}
The log likelihood of a point process with conditional intensity $\lambda$ observed on $(a,b]$ is
$\sum\log\lambda(u_j)-\int_a^b\lambda$ (Daley and Vere-Jones 2003, Prop.~7.2.III); the integral of the exponential
kernel of event $k$ over $(u_k,b]$ is $\eta(1-e^{-\beta(b-u_k)})$, and the excitation sum satisfies the recursion
(Ozaki 1979). In $\mu$, $\ell_i$ is a sum of logarithms of affine functions minus a linear term, so it is strictly
concave; its derivative $g(\mu)=\sum_j(\mu+e_j)^{-1}-S_i$ is strictly decreasing, so the stationary point is unique
when $g(0)>0$ and the constrained maximum is at $0$ otherwise. By Jensen's inequality
$\sum_j(\mu+e_j)^{-1}\ge\bar n^2/(\bar n\mu+\sum e_j)$, so $g\ge0$ at $\mu=\bar n/S_i-\bar e$: the start is at or below
the root. $h(x)=\log\sum_j(e^x+e_j)^{-1}-\log S_i$ is decreasing and smooth, and the clipped Newton iteration
converges (\texttt{test\_profile\_mu\_solves\_the\_score\_equation}). The segment parameters
$\theta_k=(\operatorname{logit}\eta,\log\beta)$ maximise the profile likelihood plus a Normal prior centred on the
pooled fit, and are drawn from the Laplace approximation (Bernstein--von Mises; van der Vaart 1998, Thm.~10.1), as
for the renewal parameters (Prop.~\ref{prop:tspost}). \texttt{test\_hawkes\_fit\_recovers\_branching\_ratio\_and\_decay}
recovers $(\eta,\beta)=(0.6,0.4)$ from simulated histories, and the compensator increments of the fit are $\mathrm{Exp}(1)$
by the time-rescaling theorem (Brown et al.\ 2002; Ogata 1988), reported as \texttt{residual\_ks} in the fit log.
\end{proof}

\begin{proposition}[Hawkes histories given their count]\label{prop:simhawkesbridge}
In the cluster representation of the Hawkes process (Hawkes and Oakes 1974), let $N=n-1$ nodes (the first event and
the $n-2$ interior events) form $r=m+1$ clusters. (a) Under Poisson($\nu$) interior immigrants, the posterior of $m$
given $N$ is proportional to $\mathrm{Pois}(m;\nu)\,\mathrm{BT}(N;m+1,\eta)$, with the Borel--Tanner law
$\mathrm{BT}(N;r,\eta)=\frac rN e^{-\eta N}(\eta N)^{N-r}/(N-r)!$. (b) Given $N$ and $r$, the offspring counts of the
$N$ nodes are $\mathrm{Multinomial}(N-r;1/N,\dots,1/N)$, independent of $\eta$, and rotating them to start after the
first minimum of their {\L}ukasiewicz walk, with the decoded trees in uniform random order, gives the law of the forest
in depth-first order. (c) Placing the other roots uniformly on the span and each child at an $\mathrm{Exp}(\beta)$
delay after its parent, truncated to the span, is exact for each child given its parent; jointly it is exact as
$\beta S\to\infty$.
\end{proposition}
\begin{proof}
(a) The total progeny of $r$ independent Galton--Watson trees with Poisson($\eta$) offspring is Borel--Tanner
(Tanner 1961; Haight and Breuer 1960); Bayes' rule with the Poisson prior gives the posterior, and
$\nu=\max(N(1-\eta)-1,0)$ matches $\mathbb E N=r/(1-\eta)$. (b) Independent Poisson counts given their sum are
multinomial with equal cells. A sequence $\xi$ with $\sum(\xi_j-1)=-r$ codes a forest of $r$ trees in depth-first
order exactly when its walk first reaches $-r$ at step $N$; by the cycle lemma (Dvoretzky and Motzkin 1947; Pitman
2006, Lemma 6.1) exactly $r$ rotations of $\xi$ do, and the rotation after the first minimum is one of them, since
after it the walk stays above its final value. The multinomial law is exchangeable, so each rotation class has the
same probability for all its members and the class determines the multiset of trees; the trees of a Galton--Watson
forest are exchangeable, so ordering them uniformly gives the forest's law. (c) Poisson immigrants given their number
are uniform; a child's delay given its parent and that it falls in the span is the truncated exponential. The joint
conditioning on all descendants falling in the span differs from sequential truncation only through clusters that
reach the end of the span, whose probability vanishes as $\beta S\to\infty$.
\texttt{test\_conditional\_bridge\_matches\_conditioned\_ogata\_histories} compares the sampler with histories
simulated by Ogata's (1981) thinning and kept when their count is $n$ (gap-share KS below 0.06 at $\beta S=40$).
\end{proof}

\begin{proposition}[Status chain]\label{prop:simchain}
With transition counts $N_k[a,b]$ and first-status counts $n^0_k$ of segment $k$, and the pooled chain
$P^0[a,\cdot]=(N[a,\cdot]+1/L)/\sum_b(N[a,b]+1/L)$, $p^0$ likewise, the posterior of the segment's chain under the
prior $P_k[a,\cdot]\sim\mathrm{Dirichlet}(c\,P^0[a,\cdot])$, $\pi_k\sim\mathrm{Dirichlet}(c\,p^0)$ is
$P_k[a,\cdot]\sim\mathrm{Dirichlet}(N_k[a,\cdot]+cP^0[a,\cdot])$, $\pi_k\sim\mathrm{Dirichlet}(n^0_k+cp^0)$,
independently over rows $a$; its mean tends to the empirical transition frequencies as the counts grow, and a segment
with no data has the pooled chain.
\end{proposition}
\begin{proof}
The likelihood of first-order Markov sequences factorises into one multinomial per origin status
(Anderson and Goodman 1957), so the Dirichlet prior is conjugate row by row (Lee, Judge and Zellner 1970). With
$c=\texttt{prior}\cdot L$ the prior carries $\texttt{prior}$ pseudo-transitions per destination on average. Segments
with fewer than \texttt{min\_entities} entities use the pooled counts, so no transition probability is read from a
handful of entities. \texttt{test\_state\_chain\_and\_holding\_shapes\_are\_recovered} recovers a three-status chain
to within 0.03.
\end{proof}

\begin{proposition}[Status-dependent gamma bridge]\label{prop:simsemimarkov}
Let the holding times after the events $j=1,\dots,n-1$ of an entity be independent
$\mathrm{Gamma}(\alpha e^{\delta_{s_j}},\beta)$ with a common scale. Given the statuses and the span $S$, the gap
shares are $\mathrm{Dirichlet}(\alpha e^{\delta_{s_1}},\dots,\alpha e^{\delta_{s_{n-1}}})$ and independent of $S$.
With $\log\alpha_i\sim N(\mu,\sigma^2)$, the marginal likelihood of entity $i$ depends on its gaps only through the
counts $C_{is}$ and log-share sums $G_{is}$ of the gaps starting in each status, and $\delta\equiv0$ is the renewal
bridge (Prop.~\ref{prop:tsbridge}).
\end{proposition}
\begin{proof}
Lukacs (1955) holds for independent gammas with different shapes and a common scale: the change of variables
$(G_1,\dots,G_m)\mapsto(S,G/S)$ factorises the density into a $\mathrm{Gamma}(\sum\alpha_j,\beta)$ density of $S$ and
the Dirichlet density of the shares. The log Dirichlet density is
$\log\Gamma(\alpha D_i)-\sum_sC_{is}\log\Gamma(\alpha e^{\delta_s})+\sum_s(\alpha e^{\delta_s}-1)G_{is}$ with
$D_i=\sum_sC_{is}e^{\delta_s}$, and the entity effect is integrated by Gauss--Hermite quadrature. $\delta$ is
identified up to a constant absorbed by $\mu$, so the most frequent status has $\delta=0$. A common scale means a
status with longer holding times is also more regular; the guide states this limit.
\end{proof}

\begin{proposition}[Scenarios leave the baseline unchanged and keep DP]\label{prop:simscenbase}
(a) With no active scenario the engine draws the same random numbers as without the \texttt{scenarios} block, so the
datasets are equal. (b) A scenario reads only the fitted parameters and the rows being sampled, never source rows,
so a dataset sampled under it from a model that is $(\varepsilon,\delta)$-DP (stage 2 of DP mode) is
$(\varepsilon,\delta)$-DP. (c) A marginal map is a deterministic function of the sampled column, so the rest of the
dataset uses the same random numbers as the baseline (common random numbers).
\end{proposition}
\begin{proof}
(a) Every hook tests for an active scenario before drawing; \texttt{test\_no\_active\_scenario\_changes\_nothing}
checks equality of all tables, and a fixed spec was checked byte for byte against the previous release. (b) Post-processing
(Dwork and Roth 2014, Prop.~2.1). (c) The map consumes no random numbers (\texttt{test\_marginal\_map\_keeps\_ranks}
compares row by row).
\end{proof}

\begin{proposition}[Tilting by rejection]\label{prop:simtilt}
Let rows $x$ of an entity table (with their segments) be drawn from the fitted model $p$ and accepted with
probability $w(x)/W$, $W\ge\sup w$. The accepted rows are i.i.d.\ from $p_w(x)=p(x)w(x)/\mathbb E_pw$, which for
$w(x)=f(x_c)$ is the distribution closest to $p$ in Kullback--Leibler divergence among those with the tilted shares of
$x_c$, and has $p_w(\cdot\mid x_c)=p(\cdot\mid x_c)$. Filling quotas $q_v$ per level $v$ of $x_c$ with the first rows
of that level in a sequence of i.i.d.\ draws gives $q_v$ i.i.d.\ rows from $p(\cdot\mid x_c=v)$.
\end{proposition}
\begin{proof}
Rejection sampling (von Neumann 1951): the accepted law has density proportional to $p\,w/W$. For a weight on $x_c$,
$p_w(x)=p(x_c)f(x_c)p(x_{-c}\mid x_c)/Z$, so the conditional is unchanged; minimising $\mathrm{KL}(q\|p)$ subject to
the marginal of $x_c$ gives $q(x)=q(x_c)p(x_{-c}\mid x_c)$, the I-projection (Csisz\'ar 1975). The quota rule is a
stopping rule per level that does not look at the other columns, so the selected rows of level $v$ are i.i.d.\ from the
conditional. Segments travel with the rows, so links of children follow the tilted segments.
\texttt{test\_target\_shares\_are\_exact\_and\_conditionals\_kept} and \texttt{test\_share\_factors\_tilt\_and\_keep\_the\_total}.
\end{proof}

\begin{proposition}[Marginal maps keep the copula]\label{prop:simmap}
Let $T$ be any map applied to the sampled values $x_1,\dots,x_n$ of a column and let $y$ be the monotone rearrangement
of $T(x)$: the sorted $T(x)$ assigned in the rank order of $x$. Then $y$ has the empirical distribution of $T(x)$ and
the same ranks as $x$, so the empirical copula of the table is unchanged; when $T$ is increasing, $y=T(x)$.
\end{proposition}
\begin{proof}
By construction $y_{(i)}=T(x)_{(i)}$ with the order of $x$, so the multiset of values is that of $T(x)$ and the ranks
are those of $x$; the empirical copula is a function of the ranks only (Deheuvels 1979; Sklar's theorem). An
increasing $T$ keeps the order, so sorting changes nothing. The rearrangement is the $L^p$-closest monotone map to $T$
(Chernozhukov, Fern\'andez-Val and Galichon 2010).
\end{proof}

\begin{proposition}[Rates and date windows by superposition and thinning]\label{prop:simrate}
Let the rows of a fact given its anchors form a Poisson process (rows per anchor, or pairs) of intensity $\lambda$.
Superposing $K$ independent draws and keeping each row with probability $c\,w(t)/K\le1$, with $t$ the row's time and
$w=f$ inside the window and $1$ outside, gives a Poisson process of intensity $c\,w\,\lambda$. For a sequence table the
calendar intensity is multiplied by $w$ on the window's cells and the count by
$g=\sum_k m_k\,\bigl(\int w\,\lambda_k/\int\lambda_k\bigr)/\sum_km_k$ (events-weighted over segments).
\end{proposition}
\begin{proof}
Superposition of independent Poisson processes adds intensities and independent thinning multiplies them by the
retention probability (Kingman 1993, Sec.~2.2 and 5.1; Lewis and Shedler 1979). For count laws that are not Poisson
(negative binomial, DP counts) the mean is multiplied by $c\,\mathbb E w$ and the dispersion changes; occupancy facts
keep unique pairs by dropping repeated pairs, so their rate is raised less than $c$ when pairs saturate (the record
reports the rows). For sequences the engine fixes the count, so the window's share of the entity's events follows
$w\lambda_k$ by the time change (Prop.~\ref{prop:tschange}) and the count follows $g$.
\texttt{test\_fact\_rate\_superposition\_and\_thinning}, \texttt{test\_intensity\_window\_on\_a\_fact\_date\_column}
and \texttt{test\_intensity\_window\_on\_a\_sequence\_table}.
\end{proof}

\begin{proposition}[Degree tempering]\label{prop:simdegree}
For degree weights $w>0$ of the rows of a segment, $w'_i=w_i^c\,\sum w/\sum w^c$ keeps the segment's total, so the
expected links of every segment (the block rates) are unchanged. For $c>1$ the shares $w'/\sum w'$ majorise $w/\sum w$
(Lorenz order: more concentrated), for $0\le c<1$ they are majorised by them, and $c=0$ gives equal shares.
\end{proposition}
\begin{proof}
The total is kept by construction. The shares $s'_i\propto s_i\cdot s_i^{c-1}$ multiply each share by a factor that
is increasing in $s_i$ when $c>1$ (decreasing when $c<1$), and reweighting by an increasing function moves mass to the
larger shares, which gives majorisation (Marshall, Olkin and Arnold 2011, Prop.~5.B.1). The tempering applies to the
heterogeneity the fit found within segments (gamma frailties or a copied table's degree posterior); when there is
none, the record says so. \texttt{test\_degree\_concentration\_tempers\_parent\_weights}.
\end{proof}
\setcounter{proposition}{\value{savedpropsim}}
